\subsection{问题三：虚拟台风情景仿真与灵敏度分析}

\subsubsection{建模思路}

第三问以问题二建立的路径--强度--环境--降水条件生成模型为基础，不另建新的降水模型。考虑到未来气候变化背景下台风路径、强度和移动速度可能偏离历史统计特征，本文将这种偏离理解为情景扰动动机，而不是确定性预报；扰动幅度仍限制在历史可比样本支持范围内。虚拟台风情景通过扰动台风路径、强度、移动速度等输入变量构造，并将扰动后的输入重新送入条件降水场生成链条，从而生成对应的半小时降水时空分布。

该过程不直接平移已有雨场，也不对基准雨场进行整体倍数放大，因此情景差异来自路径、强度、运动和环境条件共同作用下的模型响应。评价指标围绕极端降水量、强降水面积、固定地理格点持续时间和影响区域变化展开。由于 GPM 降水场的雨强单位为 $\mathrm{mm\,h^{-1}}$，而每个时次对应半小时，事件累计降水必须乘以 $0.5$ 的时间因子，不能将雨强直接累加为累计降水。

\subsubsection{虚拟台风情景构造}

本文以 KONG-REY 为基准台风构造 S0 情景，并在历史可比样本约束下设置五类扰动情景。S0 为基准情景，保持问题二输入不变；S1 为强度增强情景，仅扰动 WND 和 PRES，并保持风速增强与气压降低的一致性；S2 为路径近岸或西移情景，主要改变经度，使路径更靠近我国沿海；S3 为登陆点或近岸段偏移情景，仅在距岸最近附近施加平滑路径扰动；S4 为移动速度减慢情景，通过拉长时间轴降低沿路径推进速度；S5 为复合高风险情景，采用中等强度增强、中等近岸偏移和轻度减速的组合。各情景构造见表~\ref{tab:problem3_scenarios}。

在变量使用上，情景构造不使用 OWD，也不使用 \verb|rain_*|、\verb|centroid_*|、\verb|anisotropy|、\verb|asym_*|、\verb|quad_*|、\verb|r50|、\verb|r80|、\verb|r90| 等由降水场反推的泄漏变量。对于 S2、S3、S4 和 S5 中涉及路径或时间轴变化的情景，扰动后重新计算移动速度、移动方向、距岸距离、海陆状态、海陆覆盖比例和地形变量，以保证输入变量之间具有物理一致性。

\subsubsection{评价指标体系}

设第 $s$ 个情景在时刻 $t$ 的模型生成雨强场为 $R_{s,t}(i,j)$，单位为 $\mathrm{mm\,h^{-1}}$。由于时间分辨率为半小时，单时次累计降水为
\begin{equation}
H_{s,t}(i,j)=0.5R_{s,t}(i,j).
\end{equation}

事件累计降水定义为
\begin{equation}
C_s(i,j)=0.5\sum_t R_{s,t}(i,j).
\end{equation}

强降水面积以雨强阈值定义。以 $10\ \mathrm{mm\,h^{-1}}$ 为例，
\begin{equation}
A_{10,s}(t)=\sum_{i,j} A_{ij}\mathbf{1}\{R_{s,t}(i,j)\ge 10\},
\end{equation}
其中 $A_{ij}$ 为经纬度格点面积，按纬度进行修正。若网格分辨率为 $0.1^\circ\times0.1^\circ$，则可近似取
\begin{equation}
A_{ij}=111.32^2\times 0.1\times 0.1\times \cos\varphi_i ,
\end{equation}
其中 $\varphi_i$ 为第 $i$ 行格点纬度。固定地理格点上的强降水持续时间定义为
\begin{equation}
D_{10,s}(i,j)=0.5\sum_t \mathbf{1}\{R_{s,t}(i,j)\ge 10\}.
\end{equation}

为比较扰动情景相对基准情景的变化，定义相对变化率为
\begin{equation}
\Delta Y_s=\frac{Y_s-Y_0}{Y_0}\times 100\%,
\end{equation}
其中 $Y_0$ 表示 S0 对应指标。对于空间落区，本文进一步计算累计雨区质心偏移距离和雨区重叠率 IoU。若 $\Omega_s$ 为情景 $s$ 中超过给定阈值的地理格点集合，则
\begin{equation}
\mathrm{IoU}_s=\frac{|\Omega_s\cap\Omega_0|_A}{|\Omega_s\cup\Omega_0|_A},
\end{equation}
其中 $|\cdot|_A$ 表示按格点面积加权后的区域面积。质心偏移用于刻画降水影响中心相对 S0 的空间位移，IoU 用于刻画影响范围与基准情景的重叠程度。

\subsubsection{情景生成质量与基准一致性检验}

问题三第二步生成的半小时雨强场 shape 为 $(2841,201,201)$，其中 S0、S1、S2、S3 分别包含 421 个半小时时次，S4 包含 647 个半小时时次，S5 包含 510 个半小时时次。经检查，生成雨场中 NaN、Inf、负值和全零场数量均为 0。由于第二步雨场为以台风中心为参考的 storm-relative 网格，第三步进一步执行地理反投影，将其映射到统一的固定经纬度网格。反投影后累计降水和持续时间数组 shape 为 $(6,391,451)$。

S0 与问题二 KONG-REY 地理基准结果的事件级一致性状态为 WARNING-but-usable。S0 的最大累计降水为 $723.48\ \mathrm{mm}$，与问题二 KONG-REY 的 $722.05\ \mathrm{mm}$ 接近；但 S0 的固定地理格点 $D_{10}$ 最大值、最大累计降水位置和最大 $A_{10}$ 与问题二结果存在差异。因此，后续情景比较统一采用问题三生成链条下的 S0 作为内部对照，而不将其表述为问题二 KONG-REY 的逐像元完全复刻。需要特别说明的是，本文的 $D_{10}$ 为固定地理格点持续时间，不是 storm-relative 结构维持时间。

\subsubsection{情景结果分析}

S0--S5 的事件级核心指标见表~\ref{tab:problem3_scenario_comparison}。正文主表仅保留最大累计降水、最大强降水面积、固定地理格点最大持续时间及其相对变化，完整事件级宽表可作为附录或支撑材料。基准情景 S0 的最大累计降水为 $723.48\ \mathrm{mm}$，最大 $A_{10}$ 为 $74508.15\ \mathrm{km^2}$，固定地理格点 $D_{10}$ 最大值为 $24.5\ \mathrm{h}$。相对 S0，S4 的最大累计降水增加 $66.98\%$，$D_{10}$ 增加 $59.18\%$，总降水体量指数增加 $49.66\%$，说明移动速度减慢在模型条件下显著增强固定地理位置上的降水滞留效应。S5 的最大累计降水增加 $30.99\%$，最大 $A_{10}$ 增加 $11.06\%$，表明复合扰动可使多个风险指标同步放大。

路径扰动情景呈现不同特征。S2 的最大累计降水增加 $10.35\%$，但最大 $A_{10}$ 下降 $3.92\%$，同时累计降水超过 $100\ \mathrm{mm}$ 区域与 S0 的 IoU 为 0.67，说明该情景主要体现降水落区重新分布，而不是强降水面积的单调增加。因此，S2 的 $A_{10}$ 下降不应解释为模型失败，而应理解为路径偏移后强降水区投影到固定地理网格的位置发生改变。S3 整体指标变化较小，最大累计降水下降 $2.48\%$，固定地理格点 $D_{10}$ 下降 $4.08\%$，主要表现为近岸影响带位置调整。S1 的最大 $A_{10}$ 增加 $6.70\%$，P99 增加 $3.98\%$，反映强度增强对强降水范围和极端尾部的抬升作用。

图~\ref{fig:problem3_s0_s5_accum} 和图~\ref{fig:problem3_s0_s5_duration10} 可用于展示 S0 与 S5 在固定地理落区上的累计降水和持续时间差异；图~\ref{fig:problem3_metric_bars} 可用于比较 S0--S5 的关键指标；图~\ref{fig:problem3_path_centroid} 可用于解释路径扰动导致的降水质心偏移。

\subsubsection{灵敏度分析与物理解释}

灵敏度结果见表~\ref{tab:problem3_sensitivity}。强度增强主要影响最大雨强、P99 和强降水面积。路径近岸或西移主要改变累计雨区质心、IoU 和局地落区，因此其影响不一定表现为所有面积指标同步增大。近岸段偏移主要影响沿海风险带的位置，适合用于讨论局地防灾对象的空间转移。移动速度减慢主要影响累计降水和固定地理位置的强降水持续时间，是本组情景中对累计风险最敏感的因子。复合高风险情景表现为最大累计降水、最大强降水面积和持续时间等多项指标同步增加，但并非所有单项指标均达到最大值。

\subsubsection{小结}

情景实验表明，在历史可比扰动范围内，移动速度减慢是累计降水和持续时间最敏感因子。S4 的最大累计降水相对 S0 增加 $66.98\%$，固定地理格点 $D_{10}$ 最大值增加 $59.18\%$。

强度增强主要扩大强降水范围并增强极端尾部。S1 的最大 $A_{10}$ 增加 $6.70\%$，P99 增加 $3.98\%$，说明强度扰动更直接影响强降水覆盖和雨强尾部。

路径和登陆点变化主要改变风险落区。S2 和 S3 的结果说明，强降水影响区可随路径扰动发生明显空间转移，因此防灾减灾中不能只关注总雨量大小，还应关注降水落区的变化。
